% Part: computability
% Chapter: recursive-functions
% Section: pr-functions-computable

\documentclass[../../../include/open-logic-section]{subfiles}

\begin{document}

\olfileid[pt]{cmp}{rec}{cmp}
\olsection{As funções recursivas primitivas são computáveis}

Suponha que uma função $h$ seja definida por recursão primitiva:
\begin{eqnarray*}
h(\vec x, 0)     & = & f(\vec x) \\
h(\vec x, y + 1) & = & g(\vec x, y, h(\vec x, y))
\end{eqnarray*}
e suponha que as funções $f$ e $g$ sejam computáveis. (Usamos $\vec x$ como
abreviação de $x_0$, \dots, $x_{k-1}$.) É claro que $h(\vec x, 0)$ pode ser
calculada, pois é simplesmente $f(\vec x)$, que supomos ser computável.
Também podemos então calcular $h(\vec x, 1)$, pois $1 = 0 + 1$ e, portanto,
$h(\vec x, 1)$ é simplesmente
\begin{align*}
 h(\vec x, 1) & = g(\vec x, 0, h(\vec x, 0)) =  g(\vec x, 0, f(\vec x)).
\intertext{Podemos prosseguir dessa maneira e calcular}
h(\vec x, 2) & = g(\vec x, 1, h(\vec x, 1)) = g(\vec x, 1, g(\vec x, 0, f(\vec x)))\\
h(\vec x, 3) & = g(\vec x, 2, h(\vec x, 2)) = g(\vec x, 2, g(\vec x, 1, g(\vec x, 0, f(\vec x))))\\
h(\vec x, 4) & = g(\vec x, 3, h(\vec x, 3)) = g(\vec x, 3, g(\vec x, 2, g(\vec x, 1, g(\vec x, 0, f(\vec x)))))\\
& \vdots
\end{align*}
Assim, para calcular $h(\vec x, y)$ em geral, calculamos sucessivamente
$h(\vec x, 0)$, $h(\vec x, 1)$, \dots, até chegarmos a $h(\vec x, y)$.

Portanto, uma definição por recursão primitiva produz uma nova função
computável se as funções $f$ e $g$ são computáveis. A composição de funções
também resulta em uma função computável se as funções $f$ e $g_i$ são
computáveis.

Como as funções básicas $\Zero$, $\Succ$ e $\Proj{n}{i}$ são computáveis, e
como a composição e a recursão primitiva produzem funções computáveis a partir
de funções computáveis, segue-se que toda função recursiva primitiva é
computável.

\end{document}
